\documentclass[11pt]{article}

\usepackage[T1]{fontenc}
\usepackage{lmodern}
\usepackage{amsmath,amssymb,amsthm}
\usepackage[margin=1in]{geometry}
\usepackage{microtype}

\newtheorem{theorem}{Theorem}
\newtheorem{lemma}[theorem]{Lemma}
\newtheorem{proposition}[theorem]{Proposition}

\newcommand{\R}{\mathbb R}
\newcommand{\Bcal}{\mathcal B}
\newcommand{\Qcal}{\mathcal Q}
\newcommand{\Ical}{\mathcal I}
\newcommand{\Ibar}{\overline{\mathcal I}}
\newcommand{\Rbar}{\overline R}

\title{A Purely Analytic Aggregate-Tail Theorem}
\author{}
\date{}

\begin{document}
\maketitle

\begin{abstract}
We assemble three independently proved analytic base-frequency inequalities
with the curvature propagation from the spherical-shell argument.  The
result replaces the aggregate interval certificate: for
\(7/51\leq\rho\leq39/50\) and \(K\geq200\), the continuous reserve is
strictly positive and the discrete aggregate implication proves the
Dirichlet P\'olya inequality.  The inclusion
\(\rho_c=1/(1+2\pi)>7/51\) then gives the exact ratio range required in the
main proof.  No interval enclosure, numerical sign test, or certificate
output is used.
\end{abstract}

\section{Definitions and analytic inputs}

For \(0\leq s\leq1\), put
\[
 G_1(s)=\frac{\sqrt{1-s^2}-s\arccos s}{\pi}
       =\frac1\pi\int_s^1\arccos t\,dt.
\]
For \(0<\rho<1\), define
\[
 \eta_\rho=G_1\!\left(\max\{\rho,\tfrac12\}\right),
 \qquad
 d_\rho=\frac{2\arcsin\rho}{\pi},
 \qquad
 b_\rho=\frac{2\pi\rho}{1-\rho}.
\]
For \(K>0\), let
\[
 \mu=\rho K,\qquad
 \Rbar=\mu+\frac12,\qquad
 S=\sqrt{\Rbar^2+b_\rho K},
\]
\begin{equation}
 \Ibar(\rho,K)
 =\frac12\left[
  \Rbar S+b_\rho K\operatorname{arsinh}
       \frac{\Rbar}{\sqrt{b_\rho K}}
  -\Rbar^2\right],
 \label{eq:Ibar}
\end{equation}
and
\begin{equation}
\begin{aligned}
 \Bcal(\rho,K)={}&
 (\mu-\tfrac12)(K\eta_\rho-1)
 +\frac{d_\rho}{2}(\mu-\tfrac32)(\mu-\tfrac12)\\
 &-(1+d_\rho)\Ibar(\rho,K)
 -\frac{8(\mu+\tfrac12)}{15\sqrt\mu}.
\end{aligned}
\label{eq:B}
\end{equation}
Finally, define the curvature floor
\begin{equation}
 F(\rho)=2\rho\eta_\rho+d_\rho\rho^2
 -\frac{3(1+d_\rho)b_\rho}{800}.
 \label{eq:F}
\end{equation}

\paragraph{The three local analytic inputs.}
The following statements have complete hand proofs in the indicated local
notes.  They are the only facts imported from those files.
\begin{center}
\begin{tabular}{l|l}
source&imported conclusion on \(7/51\leq\rho\leq39/50\)\\ \hline
\texttt{B-positive.tex}
  &\(\Bcal(\rho,200)>60\)\\
\texttt{BK-positive.tex}
  &\(\Bcal_K(\rho,200)>0\)\\
\texttt{F-positive.tex}
  &\(F(\rho)>1549/84150>0\)
\end{tabular}
\end{center}
Each source uses only elementary inequalities and finite exact-rational
tables.  In particular, none of the three conclusions depends on an
interval certificate.

\paragraph{The discrete aggregate input from the main analytic argument.}
Let
\[
 A_{\rho,1}=\{x\in\R^3:\rho<|x|<1\},
\]
and let \(N_D(A_{\rho,1},K^2)\) count its Dirichlet eigenvalues strictly
below \(K^2\).
When \(\mu>1/2\), write uniquely
\[
 \mu-\frac12=J+\tau,
 \qquad J\in\mathbb N_0,\qquad 0\leq\tau<1,
\]
and set
\[
 R=\begin{cases}J,&\tau=0,\\J+1,&\tau>0.\end{cases}
 \qquad
 C=b_\rho K.
\]
For \(C>0\), define
\[
 \Ical(C,R)=\frac12\left[
 R\sqrt{R^2+C}
 +C\operatorname{arsinh}\frac{R}{\sqrt C}
 -R^2\right]
\]
and
\begin{equation}
 \Qcal(\rho,K)
 =R\lfloor K\eta_\rho\rfloor
 +d_\rho\frac{J(J+1)}2
 -(1+d_\rho)\Ical(C,R)
 -\frac{8R\tau^{5/2}}{15\sqrt\mu}.
 \label{eq:Q}
\end{equation}
The already-proved discrete aggregate lemma is the analytic implication
\begin{equation}
 \Qcal(\rho,K)\geq0
 \quad\Longrightarrow\quad
 N_D(A_{\rho,1},K^2)
 <W(\rho,K):=\frac{2(1-\rho^3)K^3}{9\pi}.
 \label{eq:discrete-implication}
\end{equation}
Its proof is the concave-head/convex-tail summation and weighted
layer-cake argument; it contains no finite computation.  We reproduce below
the short continuous-to-discrete comparison needed to invoke it.

\begin{lemma}[Continuous reserve implies the discrete reserve]
\label{lem:B-to-Q}
If \(\mu>3/2\), \(K\eta_\rho>1\), and
\(\Bcal(\rho,K)\geq0\), then
\[
 \Qcal(\rho,K)>\Bcal(\rho,K)\geq0.
\]
\end{lemma}

\begin{proof}
Put \(y=K\eta_\rho>1\).  The strict floor inequality gives
\[
 R\lfloor y\rfloor>R(y-1)
 \geq(\mu-\tfrac12)(y-1).
\]
The definition of \(J,\tau,R\) also gives
\[
 \frac{d_\rho J(J+1)}2
 >\frac{d_\rho}{2}
   (\mu-\tfrac32)(\mu-\tfrac12).
\]
Indeed, with \(x=\mu-1/2=J+\tau>1\), the function
\(x\mapsto x(x-1)\) is increasing.  If \(\tau>0\), then
\(x(x-1)<J(J+1)\); if \(\tau=0\), the difference is \(2J>0\).
The function \(x\mapsto\sqrt{x^2+C}-x\) is positive, and
\[
 \Ical(C,R)=\int_0^R(\sqrt{x^2+C}-x)\,dx.
\]
Since \(R<\mu+1/2=\Rbar\),
\[
 \Ical(C,R)<\Ical(C,\Rbar)=\Ibar(\rho,K).
\]
Finally \(R\tau^{5/2}<\Rbar\).  Substitution of these four inequalities
in \eqref{eq:Q} gives exactly
\(\Qcal(\rho,K)>\Bcal(\rho,K)\).
\end{proof}

\section{Derivative identities and curvature propagation}

For fixed \(\rho\), abbreviate \(I(K)=\Ibar(\rho,K)\) and
\[
 P=S-\Rbar.
\]
Differentiating \eqref{eq:Ibar} gives
\begin{align}
 I_K&=\rho P+\frac{b_\rho}{2}
       \operatorname{arsinh}\frac{\Rbar}{\sqrt{b_\rho K}},
 \label{eq:IK}\\
 I_{KK}
 &=\rho^2\left(\frac{\Rbar}{S}-1\right)
   +\frac{\rho b_\rho}{2S}
   +\frac{b_\rho(2\rho K-1)}{8KS}.
 \label{eq:IKK}
\end{align}
Correspondingly,
\begin{align}
 \Bcal_K={}&
 \rho(K\eta_\rho-1)+(\mu-\tfrac12)\eta_\rho
 +d_\rho\rho(\mu-1)
 -(1+d_\rho)I_K
 -\frac{2\rho(2\mu-1)}{15\mu^{3/2}},
 \label{eq:BK}\\
 E_{KK}&=\frac{\rho^2(2\mu-3)}{15\mu^{5/2}},
 \label{eq:EKK}\\
 \Bcal_{KK}&=
 2\rho\eta_\rho+d_\rho\rho^2
 -(1+d_\rho)I_{KK}+E_{KK}.
 \label{eq:BKK}
\end{align}
These are identities, not estimates.

\begin{lemma}[Exact validity conditions on the rational superset]
\label{lem:validity}
If
\[
 \frac7{51}\leq\rho\leq\frac{39}{50},
 \qquad K\geq200,
\]
then
\[
 \mu>\frac32,\qquad K\eta_\rho>1.
\]
\end{lemma}

\begin{proof}
First,
\[
 \mu=\rho K\geq\frac{1400}{51}>\frac32.
\]
The function \(G_1\) is decreasing because
\(G_1'(s)=-\arccos(s)/\pi\).  Since
\(\max\{\rho,1/2\}\leq39/50\),
\[
 \eta_\rho\geq G_1(39/50).
\]
The elementary inequality
\(\arccos t\geq\sqrt{2(1-t)}\) follows on writing
\(t=\cos\theta\), since
\(\sqrt{2(1-t)}=2\sin(\theta/2)\leq\theta\).
Combining this with \(\pi<22/7\) gives
\[
 200\eta_\rho
 \geq200G_1(39/50)
 >\frac{1400\sqrt2}{33}\left(\frac{11}{50}\right)^{3/2}.
\]
The square of the final positive quantity is
\[
 \frac{8624}{225}>1.
\]
Thus \(200\eta_\rho>1\), and \(K\eta_\rho\geq200\eta_\rho>1\).
\end{proof}

\begin{proposition}[Analytic curvature propagation]
\label{prop:curvature}
For
\[
 \frac7{51}\leq\rho\leq\frac{39}{50},
 \qquad K\geq200,
\]
one has
\begin{equation}
 \Bcal_{KK}(\rho,K)>F(\rho)>
 \frac{1549}{84150}>0.
 \label{eq:BKK-positive}
\end{equation}
\end{proposition}

\begin{proof}
Here
\[
 S>\Rbar>\rho K.
\]
In \eqref{eq:IKK}, the first term is negative.  The remaining two terms
satisfy
\begin{align*}
 \frac{\rho b_\rho}{2S}
 +\frac{b_\rho(2\rho K-1)}{8KS}
 &<\frac{\rho b_\rho}{2S}
   +\frac{2\rho Kb_\rho}{8KS}\\
 &=\frac{3\rho b_\rho}{4S}
 <\frac{3b_\rho}{4K}
 \leq\frac{3b_\rho}{800}.
\end{align*}
Therefore
\[
 I_{KK}<\frac{3b_\rho}{800}.
\]
Lemma~\ref{lem:validity} gives \(2\mu-3>0\), so
\(E_{KK}>0\).  Substitution in \eqref{eq:BKK} yields
\[
 \Bcal_{KK}
 >2\rho\eta_\rho+d_\rho\rho^2
  -\frac{3(1+d_\rho)b_\rho}{800}
 =F(\rho).
\]
The final strict rational bound is precisely the imported conclusion of
\texttt{F-positive.tex}.
\end{proof}

\section{The analytic aggregate-tail theorem}

\begin{theorem}[Continuous aggregate reserve for all \(K\geq200\)]
\label{thm:B-tail}
For every
\[
 \frac7{51}\leq\rho\leq\frac{39}{50},
 \qquad K\geq200,
\]
one has the quantitative estimate
\begin{equation}
 \Bcal(\rho,K)>
 60+\frac{1549}{168300}(K-200)^2>0.
 \label{eq:B-tail-lower}
\end{equation}
\end{theorem}

\begin{proof}
Fix \(\rho\) in the stated interval.  The imported local lemmas give
\[
 \Bcal(\rho,200)>60,\qquad
 \Bcal_K(\rho,200)>0.
\]
Proposition~\ref{prop:curvature} gives
\[
 \Bcal_{KK}(\rho,s)>\frac{1549}{84150}
 \qquad(s\geq200).
\]
Twice integrating this strict inequality from \(200\) to \(K\) gives
\begin{align*}
 \Bcal(\rho,K)
 &=\Bcal(\rho,200)
 +(K-200)\Bcal_K(\rho,200)
 +\int_{200}^K(K-s)\Bcal_{KK}(\rho,s)\,ds\\
 &>60+\frac12\frac{1549}{84150}(K-200)^2,
\end{align*}
including \(K=200\) by the strict imported base value.  This is
\eqref{eq:B-tail-lower}.
\end{proof}

\begin{theorem}[Purely analytic aggregate tail for spherical shells]
\label{thm:aggregate-tail}
Let
\[
 \rho_c=\frac1{1+2\pi}.
\]
For every
\[
 \rho_c\leq\rho\leq\frac{39}{50},
 \qquad K\geq200,
\]
the Dirichlet counting function of the unit outer-radius spherical shell
satisfies
\[
 N_D(A_{\rho,1},K^2)
 <\frac{2(1-\rho^3)K^3}{9\pi}.
\]
\end{theorem}

\begin{proof}
The exact Archimedean bound \(\pi<22/7\) gives
\[
 1+2\pi<1+\frac{44}{7}=\frac{51}{7},
 \qquad\text{hence}\qquad
 \rho_c>\frac7{51}.
\]
(Also \(\pi>3\) gives \(\rho_c<1/7<39/50\).)
Thus the stated ratio interval is contained in the rational superset used
in Theorem~\ref{thm:B-tail}.

On this superset, Lemma~\ref{lem:validity} verifies
\(\mu>3/2\) and \(K\eta_\rho>1\), while
Theorem~\ref{thm:B-tail} gives \(\Bcal(\rho,K)>0\).
Lemma~\ref{lem:B-to-Q} therefore gives
\(\Qcal(\rho,K)>0\).  The analytic discrete implication
\eqref{eq:discrete-implication} now proves the asserted strict P\'olya
bound.
\end{proof}

\paragraph{Dependency boundary.}
The only local-file imports are the three base inequalities listed at the
start.  The derivative identities, curvature propagation, continuous-to-
discrete comparison, rational-domain checks, and two integrations are all
displayed in this note.  The remaining discrete implication is the
previously proved analytic layer-cake lemma from the main manuscript.
No aggregate certificate predicate or interval-arithmetic result appears in
the dependency chain.

\end{document}
